\section{Benchmark 设计}

当前 benchmark 的第一原则是统一数据边界。即使同属 envelope reconstruction，不同项目对 train/val/test 的定义、窗口长度、滑窗步长、recording 聚合方式以及 subject conditioning 的处理都可能不同。本项目在本轮仅要求完成最小真实闭环，因此使用一个小型公开 sample 作为 pipeline validation 的载体，但输出格式、打分脚本和边界检查方式已经按后续多数据集扩展的标准设计。

第二原则是统一输出 schema。每个模型都必须最终对接相同的输出层，包括 prediction 样本、recording 级指标、subject 级指标与 run manifest。这样后续无论是 classical baseline、simple neural baseline 还是 faithful port，都可以被同一套归档、可视化与统计工具消费。

第三原则是明确任务矩阵。当前主任务是 envelope reconstruction；match-mismatch 和 AAD 暂作为派生或扩展任务。对于后续数据集接入，必须先标记每个数据是否适合 reconstruction、是否存在 competing-speaker 条件、是否支持 AAD 或 match-mismatch 扩展。

\begin{table}[ht]
\centering
\caption{当前数据集与任务矩阵（以评估角色而非论文结论为准）}
\begin{tabular}{llll}
\toprule
数据集 & 当前状态 & 支持任务 & 当前角色 \\
\midrule
hugo\_sample\_tf64\_p00 & 已接入并验证 & reconstruction & pipeline validation \\
hugo\_sample\_tf64 & 已存在但本轮未重跑 & reconstruction & development / smoke set \\
weissbart\_tf64\_p00 & 本轮已运行 & reconstruction & article-grade dataset pilot \\
etard\_tf64\_p00 & 本轮已运行 & reconstruction & article-grade dataset pilot \\
weissbart\_tf64 & 已处理，待全被试统一重跑 & reconstruction & main benchmark candidate \\
etard\_tf64 & 已作为 P00 pilot 入口验证，待全被试统一重跑 & reconstruction & main benchmark candidate \\
SparrKULee / KULeuven AAD & 计划接入 & reconstruction / AAD & benchmark expansion \\
Fuglsang / DTU / 其它外部集 & 部分可得 & generalization / AAD & external validation \\
\bottomrule
\end{tabular}
\end{table}


\begin{table}[ht]
\centering
\caption{当前统一指标与结果 schema}
\begin{tabular}{lll}
\toprule
层级 & 文件 / 指标 & 说明 \\
\midrule
sample 级 & prediction\_samples.csv & 连续预测与 target 的统一导出 \\
recording 级 & recording\_metrics.csv & 主指标为 Pearson correlation \\
subject 级 & subject\_metrics.csv & recording 指标的 subject 聚合 \\
run 级 & run\_manifest.json & 记录模型、配置、artifact scope 与命令来源 \\
附加指标 & canonical / mm accuracy & 仅用于特定模型的附加说明，不并入主比较列 \\
\bottomrule
\end{tabular}
\end{table}

